% Apartados propuestos y aún no redactados. Al redactar uno, se mueve a su capítulo sin
% etiqueta y se elimina de este archivo.
% Uso: \pendiente{número recomendado}{título}{ubicación}{etiqueta}{contenido previsto}
\newcommand{\pendiente}[5]{%
	\addcontentsline{toc}{section}{#1 #2 (#3) [#4]}%
	\noindent\textbf{#1 #2} [#4]\\
	\textit{Ubicación:} #3.\\
	#5\par\medskip
}

\pendiente{2.4}{Optimism}{entre Aztec y Puentes}{Necesario - Importante}{%
	Optimism y el OP Stack como rollup optimista de Ethereum. Estructura de bloques y
	compatibilidad con la EVM, incluida la equivalencia de gas y de contratos precompilados.
	Pruebas de fraude, período de impugnación y finalidad. Mensajería nativa con Ethereum y su
	diferencia con el puente propuesto. Optimism Sepolia como cadena de destino en la que se
	despliega el contrato verificador.}
